\chapter{Selection of the plateau and correlation aging}\label{sec:stationary-selection}

In this section, we prove the parts of \cref{thm:main-aging} that concern the plateau, the energy, and the loss of memory. We begin with \cref{sel:positive-marginality}, which states that every positive plateau of a stationary cluster is marginal. In \cref{sel:all-temperature}, we prove that the translated kernels, the multiplier, and the energy converge at every temperature. \Cref{sel:selected-sector} identifies the plateau and the energy for $\beta\ge\bd$, and \cref{sel:fixed-wait-forgetting} shows that the solution forgets every fixed earlier state. \Cref{sel:aging} states the two ordered limits of the correlation. At the end of \cref{sec:aging}, we combine these results into a proof of \cref{thm:main-aging}.

We take as input two facts about prefix centers, which we state as \cref{prop:threshold-inputs} in \cref{sec:threshold-inputs} and prove in \cref{sec:tree-ceilings,app:general-selection}. Besides these two facts, we use the stationary clusters of \cref{prop:stationary-clusters}, the scalar constraints of \cref{prop:scalar-clusters}, the uniqueness result \cref{lem:stationary-uniqueness}, the causal tree realization of \cref{prop:geometry-tree-realization}, the reflection identity of \cref{sel:prefix-reflection}, and the uniform bounds of \cref{sec:ck-structure}. The convergence proved in \cref{sel:all-temperature} and the marginality proved in \cref{sel:positive-marginality} are exactly the hypotheses of part~\ref{item:threshold-forgetting} of \cref{prop:threshold-inputs}. We use the selected profile, the marginality relation, and the forgetting theorem throughout \cref{part:response,part:geometry,part:relaxation}.

\paragraph{Standing assumptions and notation.}
Throughout this section, we fix $p\ge3$ and $\beta>0$, and $(C,R,\lambda)$ denotes the uniformly initialized hard CK solution of \cref{sec:model}. All long-time limits below concern this deterministic solution, for which the dimension limit has already been taken on each fixed time horizon. Recall from \eqref{eq:energy-identity} that the energy $e_\beta(t)$ satisfies $\lambda(t)=1+p\beta e_\beta(t)$. For $T\ge0$, the translated kernels are $C_T(a,b)=C(T+a,T+b)$ for $a,b\ge-T$, $R_T(a,b)=R(T+a,T+b)$ for $-T\le b\le a$, and $\lambda_T(a)=\lambda(T+a)$ for $a\ge-T$.

We fix the Fisher majorant $a_F$ and the dimension subsequence $(N_j)$ of \cref{thm:fisher-clock} once, before choosing any waiting-time sequence or stationary cluster. All comparisons with finite-dimensional forks use this same dimension subsequence. We may pass to subsequences of waiting times to obtain stationary clusters, without changing the dimension extraction. The majorant is an auxiliary choice; the CK solution and the stationary profile identified below are unique.

As in \cref{prop:stationary-clusters}, a \emph{stationary cluster} is a triple $(c,r,\ell)$ obtained as the locally uniform limit of $(C_{T_n},R_{T_n},\lambda_{T_n})$ along some sequence $T_n\to\infty$. \Cref{prop:stationary-clusters} implies that $c(0)=r(0)=1$, $c'=-r$, $0\le c\le1$, and $r\ge0$, that the plateau $q=\lim_{\tau\to\infty}c(\tau)$ lies in $[0,1)$, and that $\int_0^\infty r(u)\,\dd u=1-q$. By \eqref{eq:response-derivative-bounds}, we have $0\le R\le C$, and passing to the limit gives $r\le c$. As in \cref{prop:stationary-clusters,sel:prefix-reflection}, we write $e_\infty=(\ell-1)/(p\beta)$ for the energy of the cluster. With $B=\beta^2p$ and $\gamma=\ell-B$, we use the functions of \cref{prop:scalar-clusters},
\[
 \phi(y)=\gamma+By^{p-1},\qquad F(y)=\phi(y)(1-y)-1.
\]
Recall from that proposition that $F(q)=0$, that $F(y)<0$ for $q<y<1$, and that
\[
 \ell=\frac1{1-q}+\beta^2p\bigl(1-q^{p-1}\bigr).
\]

We also use the tree realization of \cref{sec:geometry-core}. The unit vectors $v_t$ satisfy $\ip{v_t}{v_s}=C(t,s)$, and $v_t$ lies in the closed subspace $\mathcal H_t$ generated by trees whose noise leaves carry times at most $t$. We write $P_T$ for the orthogonal projection onto $\mathcal H_T$. The tree energy is $h(v)=\sqrt p\,\ip{v}{\mathcal I(v^{\otimes m})}$ with $m=p-1$, and $h(v_t)=e_\beta(t)$.

\section{Prefix centers and the threshold inputs}\label{sec:threshold-inputs}

For $T\ge0$ and $a\ge0$, we define the prefix center
\[
 M_T(a)=P_Tv_{T+a}.
\]
Let $T_n\to\infty$ be a sequence along which the translated kernels converge to a stationary cluster $(c,r,\ell)$ with energy $e_\infty$. Recall from \cref{sel:prefix-reflection} that, for every fixed $L\ge0$,
\[
\begin{aligned}
 \norm{M_{T_n}(L)}^2&\longrightarrow d_L=c(2L),\\
 h\bigl(M_{T_n}(L)\bigr)&\longrightarrow
 h_L=e_\infty+\beta p\int_0^Lr(\tau)\bigl\{c(2L-\tau)^{p-1}-c(\tau)^{p-1}\bigr\}\,\dd\tau.
\end{aligned}
\]
Note that $h$ is homogeneous of degree $p$. If $d_L>0$, the energy of the normalized center $M_{T_n}(L)/\norm{M_{T_n}(L)}$ converges to $h_L/d_L^{p/2}$.

We now explain why the prefix centers are the right objects for selecting the plateau. The state $v_{T+L}$ is built from the initial vector and from the noise leaves at times up to $T+L$. The projection $P_T$ keeps the part built from the initial vector and noise received by time $T$, including memory interactions after $T$. In physical terms, we may think of $M_T(L)$ as the average of the state at time $T+L$ over independent continuations of the thermal noise after time $T$. %
This interpretation goes back to the trajectories with a common past and independent futures studied by Cugliandolo and Dean \cite[arXiv version~1, Section~II]{cugliandolodean1995} and by Barrat, Burioni, and M\'ezard \cite[Section~4]{bbm1996}, and to the averages over cloned futures proposed by Franz and Virasoro \cite[arXiv version~3, Section~III]{franzvirasoro2000}. In agreement with this interpretation, the limiting squared norm $d_L=c(2L)$ is the stationary correlation at the reflected lag $2L$, which is the overlap of two continuations of length $L$ from a common past; this reflected formula appears in \cite[Section~2, Equation~(9)]{bbm1996}. Time-averaged magnetizations of a similar kind appear in the analysis of Cugliandolo and Kurchan \cite[arXiv version~1, Section~5.1]{ck1994}, and noise-averaged magnetizations and their relation to the threshold appear in the dynamical TAP approach of Biroli \cite[arXiv version~2, Section~III.B]{biroli1999}. We use three properties of the centers.
\begin{itemize}
\item Every earlier state $v_s$ with $s\le T$ lies in $\mathcal H_T$, so that $\ip{v_{T+L}}{v_s}=\ip{M_T(L)}{v_s}$. The correlation of a late state with a fixed earlier state is carried entirely by the center.
\item For a stationary cluster with $q>0$, the normalized limiting energies of the centers satisfy $h_L/d_L^{p/2}\le\Eth$ for every $L\ge0$. Here $\Eth=2\sqrt{(p-1)/p}$ is the threshold energy of the pure $p$-spin landscape \cite[Equation~(2.13)]{abc2013}.
\item Suppose that the whole translated family converges to one stationary profile whose plateau is positive and marginal. Then the normalized energies of the centers approach $\Eth$ as $L\to\infty$, and this forces the solution to forget every fixed earlier state.
\end{itemize}
The first property is exact. The ceiling in the second and the forgetting in the third are the content of the following proposition, which is the only input that we take from \cref{sec:tree-ceilings,app:general-selection}.

\begin{proposition}[Threshold inputs]\label{prop:threshold-inputs}
Fix $p\ge3$ and $\beta>0$.
\begin{enumerate}[label=(\alph*)]
\item\label{item:threshold-ceiling}
Let $T_n\to\infty$ be a sequence along which the translated kernels converge to a stationary cluster $(c,r,\ell)$ with plateau $q>0$, as in \cref{prop:stationary-clusters}, and let $d_L$ and $h_L$ be the limits in \eqref{sel:prefix-norm} and \eqref{sel:prefix-energy}. Then, for every $L\ge0$,
\begin{equation}\label{eq:threshold-input-ceiling}
 h_L\le \Eth\, d_L^{p/2}.
\end{equation}
\item\label{item:threshold-forgetting}
Suppose that the translated kernels converge along every sequence $T\to\infty$ to one stationary profile $(c,r)$, and that its plateau $q$ is positive and satisfies $\beta^2p(p-1)q^{p-2}(1-q)^2=1$. Then, for every $0\le S<\infty$,
\begin{equation}\label{eq:threshold-input-forgetting}
 \lim_{t\to\infty}\sup_{0\le s\le S}|C(t,s)|=0.
\end{equation}
\end{enumerate}
For $p\in\{3,4\}$, both parts are proved in \cref{sec:tree-ceilings}. For every $p\ge3$, they are proved in \cref{app:general-selection}.
\end{proposition}

For $p\in\{3,4\}$, we prove both parts in \cref{sec:tree-ceilings} using a deterministic inequality on the tree Hilbert space (\cref{thm:geometry-tree-ceiling,thm:geometry-quartic-tree-ceiling}). It states that a unit tree vector has energy strictly below $\Eth$, and that near equality forces its mass to leave every fixed set of tree heights. Applying this inequality to the normalized centers, we obtain part~\ref{item:threshold-ceiling}. For part~\ref{item:threshold-forgetting}, we show that the mass of the centers escapes from every fixed height as they approach the threshold, while an earlier state $v_s$ with $s\le S$ is approximated uniformly by finitely many heights. From the same argument, we also obtain the quantitative estimate \cref{sel:threshold-centers}, which we use again in \cref{resp:chapter,sec:geometry-results,sec:three-time}.

For every $p\ge3$, we prove both parts in \cref{app:general-selection}. There we replace the projection by the mean of finitely many independent futures of the finite-dimensional dynamics from a common past. Using the fixed-horizon approximation of Sellke \cite[Lemmas~3.1--3.2]{sellke2024}, we approximate the normalized mean by globally Lipschitz functions of the disorder coefficients. The algorithmic threshold of Huang and Sellke \cite[Theorem~1 and Corollary~1.8]{hs2023} bounds the energy of these approximations by $\Eth$ in the dimension limit, and we transfer this bound to the mean by sending the approximation error to zero. Part~\ref{item:threshold-ceiling} follows from \cref{sel:finite-forks} by letting the number of futures tend to infinity at fixed $L$. Part~\ref{item:threshold-forgetting} follows from \cref{sel:general-fixed-wait} and the finite-net argument in \cref{sec:general-proof-inputs}. Its proof combines the external-field form of the same threshold with adaptive Gaussian conditioning and an Euler approximation of the earlier state.

The references to later sections do not create a circular argument. The proofs of part~\ref{item:threshold-ceiling} in \cref{sec:tree-ceilings,app:general-selection} apply to an arbitrary stationary cluster with positive plateau, and they use only results of \cref{sec:model,sec:ck-structure,sec:entropy-clusters,sec:scalar,sec:geometry-core}. Part~\ref{item:threshold-forgetting} contains in its hypotheses the two facts that we deduce from part~\ref{item:threshold-ceiling} in this section: convergence of the whole translated family, and marginality of the plateau. Besides these hypotheses, its proofs use the algebraic identity \eqref{sel:center-amgm}, which we prove in \cref{lem:center-energy} without using \cref{prop:threshold-inputs}. In this section, the order of deduction is as follows:
\begin{itemize}
\item part~\ref{item:threshold-ceiling} and \cref{lem:center-energy} imply marginality, \cref{sel:positive-marginality};
\item marginality implies convergence at every temperature, \cref{sel:all-temperature}, and the selection of the plateau, \cref{sel:selected-sector};
\item these results verify the hypotheses of part~\ref{item:threshold-forgetting}, which then gives forgetting, \cref{sel:fixed-wait-forgetting}, and correlation aging, \cref{sel:aging}.
\end{itemize}

\section{Marginality of positive plateaus}\label{sec:marginality}

We set
\begin{equation}
 q_{\mathrm m}=\frac{p-2}{p},\qquad
 \beta_{\mathrm m}^{\,2}=\frac{p^{p-1}}{4(p-1)(p-2)^{p-2}},\qquad
 \Eth=2\sqrt{\frac{p-1}{p}}.
 \label{sel:spinodal-constants}
\end{equation}
To explain the role of $q_{\mathrm m}$ and $\beta_{\mathrm m}$, let $G(q)=p(p-1)q^{p-2}(1-q)^2$ for $q\in[0,1]$, so that the marginality relation $\beta^2p(p-1)q^{p-2}(1-q)^2=1$ reads $\beta^2G(q)=1$. We have
\[
 G'(q)=p(p-1)q^{p-3}(1-q)\bigl\{(p-2)-pq\bigr\}.
\]
The function $G$ vanishes at $0$ and $1$, increases strictly on $[0,q_{\mathrm m}]$, and decreases strictly on $[q_{\mathrm m},1]$. Its maximum is
\[
 G(q_{\mathrm m})=\frac{4(p-1)(p-2)^{p-2}}{p^{p-1}}=\beta_{\mathrm m}^{-2}.
\]
In other words, $\beta_{\mathrm m}$ is the inverse temperature at which $\beta^2G$ first reaches one on $[0,1]$. For $\beta\ge\beta_{\mathrm m}$, we denote by $q_+(\beta)$ the unique root of $\beta^2G(q)=1$ in $[q_{\mathrm m},1)$; at $\beta=\beta_{\mathrm m}$, it is the double root $q_{\mathrm m}$. For $\beta>\beta_{\mathrm m}$, the equation has a second root, which lies in $(0,q_{\mathrm m})$, and for $\beta<\beta_{\mathrm m}$ it has no root in $[0,1]$.

The next lemma computes the limit of the center energies of \cref{sel:prefix-reflection} as $L$ tends to infinity, and compares it with $\Eth$. It is a computation with the scalar constraints of \cref{prop:scalar-clusters} and uses no threshold input. We use it again in \cref{sec:tree-ceilings,app:general-selection}.

\begin{lemma}[Limiting energy of prefix centers]\label{lem:center-energy}
Let $(c,r,\ell)$ be a stationary cluster with plateau $q$, as in \cref{prop:stationary-clusters}, let $e_\infty=(\ell-1)/(p\beta)$, and let $d_L$ and $h_L$ be defined by \eqref{sel:prefix-norm} and \eqref{sel:prefix-energy}. Then $d_L\to q$ and $h_L\to h$ as $L\to\infty$, where
\begin{equation}
 h=e_\infty-\beta\bigl\{1-q^p-pq^{p-1}(1-q)\bigr\}.
 \label{sel:limiting-center-energy}
\end{equation}
If $q>0$, define $z=\beta\sqrt{p(p-1)}\,q^{(p-2)/2}(1-q)>0$. Then
\begin{equation}
 \frac{h}{q^{p/2}}
 =\sqrt{\frac{p-1}{p}}\left(z+\frac1z\right)
 \ge \Eth,
 \label{sel:center-amgm}
\end{equation}
and equality holds in \eqref{sel:center-amgm} if and only if $\beta^2p(p-1)q^{p-2}(1-q)^2=1$.
\end{lemma}

\begin{proof}
As $c$ is nonincreasing with limit $q$, we have $d_L=c(2L)\to q$. For $h_L$, we extend the integrand in \eqref{sel:prefix-energy} by zero outside $[0,L]$ and write the integral as an integral over $[0,\infty)$ of
\[
 g_L(\tau)=
 \begin{cases}
 r(\tau)\bigl\{c(2L-\tau)^{p-1}-c(\tau)^{p-1}\bigr\},&0\le\tau\le L,\\
 0,&\tau>L.
 \end{cases}
\]
For fixed $\tau$, we have $2L-\tau\to\infty$ as $L\to\infty$, and $g_L(\tau)$ converges to $r(\tau)\{q^{p-1}-c(\tau)^{p-1}\}$. Since $0\le c\le1$, we have $|g_L(\tau)|\le r(\tau)$, and $r$ is integrable with $\int_0^\infty r=1-q$. By dominated convergence, we obtain
\[
 \lim_{L\to\infty}h_L
 =e_\infty+\beta pq^{p-1}(1-q)-\beta p\int_0^\infty r(\tau)c(\tau)^{p-1}\,\dd\tau.
\]
Using $r=-c'$, $c(0)=1$, and $c(\tau)\to q$, we find
\[
 \int_0^\infty r(\tau)c(\tau)^{p-1}\,\dd\tau
 =-\frac1p\int_0^\infty\bigl(c^p\bigr)'(\tau)\,\dd\tau
 =\frac{1-q^p}p.
\]
Inserting this into the previous display, we obtain \eqref{sel:limiting-center-energy}.

Now suppose that $q>0$. Using the formula for $\ell$ in \eqref{eq:plateau-multiplier}, we write the energy of the cluster as
\[
 e_\infty=\frac{\ell-1}{p\beta}=\frac{q}{p\beta(1-q)}+\beta\bigl(1-q^{p-1}\bigr).
\]
We insert this into \eqref{sel:limiting-center-energy} and note that $(1-q^{p-1})-(1-q^p)=-q^{p-1}(1-q)$. This gives
\[
 h=\frac{q}{p\beta(1-q)}+\beta(p-1)q^{p-1}(1-q).
\]
By the definition of $z$, we have
\[
 \sqrt{\frac{p-1}p}\,z=\beta(p-1)q^{(p-2)/2}(1-q),
 \qquad
 \sqrt{\frac{p-1}p}\,\frac1z=\frac{q^{-(p-2)/2}}{p\beta(1-q)},
\]
and the sum of the two right sides is $h/q^{p/2}$. This proves the identity in \eqref{sel:center-amgm}. For $z>0$, we have $z+z^{-1}\ge2$, with equality if and only if $z=1$. Recalling that $2\sqrt{(p-1)/p}=\Eth$, we deduce the inequality in \eqref{sel:center-amgm}, and equality holds if and only if $z=1$, that is, if and only if $z^2=\beta^2p(p-1)q^{p-2}(1-q)^2=1$.
\end{proof}

\begin{lemma}[Marginality of every positive cluster plateau]\label{sel:positive-marginality}
Let $(c,r,\ell)$ be a stationary cluster, as in \cref{prop:stationary-clusters}, with plateau $q>0$. Then
\begin{equation}
 \beta^2p(p-1)q^{p-2}(1-q)^2=1,
 \qquad q\ge q_{\mathrm m}.
 \label{sel:positive-plateau-equation}
\end{equation}
\end{lemma}

\begin{proof}
First, part~\ref{item:threshold-ceiling} of \cref{prop:threshold-inputs} gives an upper bound on the limiting center energy that matches the lower bound in \eqref{sel:center-amgm}, and this forces marginality. Second, the strict sign condition of \cref{prop:scalar-clusters} excludes the smaller root.

\paragraph{Step 1: marginality.}
Let $T_n\to\infty$ be a sequence along which the translated kernels converge to $(c,r,\ell)$. Since $q>0$, part~\ref{item:threshold-ceiling} of \cref{prop:threshold-inputs} applies along this sequence, and \eqref{eq:threshold-input-ceiling} states that $h_L\le\Eth d_L^{p/2}$ for every $L\ge0$. Letting $L\to\infty$ and using $d_L\to q$ and $h_L\to h$ from \cref{lem:center-energy}, we obtain
\[
 h\le\Eth\,q^{p/2}.
\]
Together with the lower bound in \eqref{sel:center-amgm}, this forces $h/q^{p/2}=\Eth$. The equality case of \cref{lem:center-energy} then implies the first relation in \eqref{sel:positive-plateau-equation}.

\paragraph{Step 2: the larger root.}
We use the functions $\phi$ and $F$ from the beginning of this section; recall that $F(q)=0$ and $F(y)<0$ for $q<y<1$. Differentiating, we have
\[
 F'(y)=B(p-1)y^{p-2}(1-y)-\phi(y),\qquad
 F''(y)=B(p-1)y^{p-3}\bigl\{(p-2)-py\bigr\}.
\]
Since $\phi(q)=1/(1-q)$ by \eqref{eq:plateau-multiplier}, we find
\[
 F'(q)=\frac{\beta^2p(p-1)q^{p-2}(1-q)^2-1}{1-q}=0,
\]
where we used Step~1 in the last equality. Suppose that $0<q<q_{\mathrm m}$. Then $F''(q)>0$. Since $F$ is a polynomial with $F(q)=F'(q)=0$, Taylor's theorem implies that
\[
 F(q+\varepsilon)=\tfrac12F''(q)\varepsilon^2+O(\varepsilon^3)>0
\]
for all sufficiently small $\varepsilon>0$. This contradicts $F(q+\varepsilon)<0$. We conclude that $q\ge q_{\mathrm m}$.
\end{proof}

The first relation in \eqref{sel:positive-plateau-equation} is the marginality condition of the asymptotic solution of Cugliandolo and Kurchan \cite[Equation~(7)]{ck1993}, written in the conventions of \cref{sec:model}. In view of the discussion after \eqref{sel:spinodal-constants}, the relations \eqref{sel:positive-plateau-equation} say that $\beta\ge\beta_{\mathrm m}$ and $q=q_+(\beta)$. A stationary cluster with positive plateau can occur only for $\beta\ge\beta_{\mathrm m}$, and its plateau is then determined by $\beta$.

\section{Convergence at every fixed temperature}\label{sec:all-temperature}

\begin{proposition}[Convergence at every fixed temperature]\label{sel:all-temperature}
For every $p\ge3$ and $\beta>0$, the multiplier $\lambda(t)$ and the energy $e_\beta(t)$ have limits $\ell_\beta$ and $e_\beta(\infty)$ as $t\to\infty$. All bounded-lag translates of $C$ and $R$ converge locally uniformly to a unique stationary FDT profile $(c,r)$, that is, for every $0\le A<\infty$,
\begin{equation}
 \sup_{0\le s\le t\le A}
 \bigl\{|C(T+t,T+s)-c(t-s)|+|R(T+t,T+s)-r(t-s)|\bigr\}
 \longrightarrow0\qquad(T\to\infty).
 \label{eq:all-temperature-window}
\end{equation}
Its plateau $q$ belongs to the following set:
\begin{equation}
 q\in\begin{cases}
 \{0\}, & 0<\beta<\beta_{\mathrm m},\\
 \{0,q_+(\beta)\}, & \beta\ge\beta_{\mathrm m}.
 \end{cases}
 \label{sel:two-possible-plateaux}
\end{equation}
The positive branch is used only on its domain $\beta\ge\beta_{\mathrm m}$. The limiting multiplier is
\[
 \ell_\beta=\frac1{1-q}+\beta^2p\bigl(1-q^{p-1}\bigr),
\]
as in \eqref{eq:plateau-multiplier}, and $e_\beta(\infty)=(\ell_\beta-1)/(p\beta)$.
\end{proposition}

\begin{proof}
First, the multiplier of every stationary cluster lies in a set with at most two elements. Second, the late limit points of $\lambda$ form an interval, each point of which is the multiplier of a stationary cluster; combined with the first step, this forces $\lambda$ to converge. Third, we identify all stationary clusters using \cref{lem:stationary-uniqueness} and upgrade subsequential convergence to convergence of the whole translated family by precompactness.

\paragraph{Step 1: possible multipliers.}
Let $(c,r,\ell)$ be a stationary cluster with plateau $q$. If $q=0$, then \eqref{eq:plateau-multiplier} gives $\ell=1+\beta^2p$. If $q>0$, then \cref{sel:positive-marginality} and the discussion after it imply that $\beta\ge\beta_{\mathrm m}$ and $q=q_+(\beta)$, and by \eqref{eq:plateau-multiplier} we have $\ell=(1-q_+(\beta))^{-1}+\beta^2p(1-q_+(\beta)^{p-1})$. In both cases, $\ell$ belongs to the set
\[
 \mathcal L_\beta=
 \begin{cases}
 \{1+\beta^2p\}, & \beta<\beta_{\mathrm m},\\[2pt]
 \bigl\{1+\beta^2p,\ (1-q_+(\beta))^{-1}+\beta^2p\bigl(1-q_+(\beta)^{p-1}\bigr)\bigr\}, & \beta\ge\beta_{\mathrm m},
 \end{cases}
\]
which has at most two elements.

\paragraph{Step 2: convergence of the multiplier.}
Recall from \cref{prop:uniform-bounds,prop:global-regularity} that $\lambda$ is continuous on $[0,\infty)$ and satisfies $1\le\lambda\le\Lambda$. Let
\[
 a=\liminf_{t\to\infty}\lambda(t),\qquad b=\limsup_{t\to\infty}\lambda(t),
\]
which are finite. We say that $y$ is a late limit point of $\lambda$ if there exists a sequence $T_n\to\infty$ such that $\lambda(T_n)\to y$. We claim that every $y\in[a,b]$ is a late limit point. For $y\in\{a,b\}$, this holds by the definition of the lower and upper limits. If $a<y<b$, then for every $n$ there exist times $t_1,t_2\ge n$ such that $\lambda(t_1)<y<\lambda(t_2)$. By the intermediate value theorem, there exists $T_n$ between $t_1$ and $t_2$ such that $\lambda(T_n)=y$, and $T_n\ge n$. This proves the claim.

Now let $y\in[a,b]$, and let $T_n\to\infty$ be such that $\lambda(T_n)\to y$. By \cref{prop:stationary-clusters}, a subsequence of the translated family $(C_{T_n},R_{T_n},\lambda_{T_n})$ converges locally uniformly to a stationary cluster $(c,r,\ell)$, and along it $\lambda_{T_n}$ converges to the constant $\ell$. Evaluating at offset zero, we find $\ell=y$, and Step~1 gives $y\in\mathcal L_\beta$. We have shown that
\[
 [a,b]\subseteq\mathcal L_\beta.
\]
An interval contained in a finite set is a single point. Hence $a=b$, and $\lambda(t)$ converges to a limit $\ell_\beta\in\mathcal L_\beta$. By \eqref{eq:energy-identity}, we have $e_\beta(t)=(\lambda(t)-1)/(p\beta)$, and the energy converges to $e_\beta(\infty)=(\ell_\beta-1)/(p\beta)$.

\paragraph{Step 3: identification of the profile.}
Let $(c,r,\ell)$ be any stationary cluster, obtained along a sequence $T_n\to\infty$. Since $\lambda(T_n+a)\to\ell_\beta$ for every offset $a$, we have $\ell=\ell_\beta$. By \eqref{eq:cluster-fdt-properties} and \eqref{eq:stationary-cluster-response}, the pair $(c,r)$ solves
\[
 c'=-r,\qquad
 r'=-\ell_\beta r+\beta^2p(p-1)\bigl(c^{p-2}r\bigr)*r,\qquad
 c(0)=r(0)=1,
\]
and it is bounded, since $0\le r\le c\le1$. \Cref{lem:stationary-uniqueness} states that this system has at most one locally bounded solution. All stationary clusters therefore coincide; we denote their common profile by $(c,r)$. The argument uses only the value $\ell_\beta$ of the multiplier. In particular, it applies when the two elements of $\mathcal L_\beta$ coincide, which happens at $\beta=\bd$. Every stationary cluster also has the same energy $e_\infty=(\ell_\beta-1)/(p\beta)=e_\beta(\infty)$.

Suppose that \eqref{eq:all-temperature-window} fails for some $0\le A<\infty$. Then there exist $\varepsilon>0$ and a sequence $T_n\to\infty$ such that, for every $n$, the supremum in \eqref{eq:all-temperature-window} at $T=T_n$ is at least $\varepsilon$. Using the precompactness in \cref{prop:stationary-clusters}, we extract a subsequence of the translates that converges locally uniformly to a stationary cluster, which is $(c,r,\ell_\beta)$ by the preceding paragraph. Along this subsequence, the supremum over the compact set $\{0\le s\le t\le A\}$ tends to zero, which is a contradiction. This proves \eqref{eq:all-temperature-window}.

Finally, if the plateau $q$ of $c$ is positive, then \cref{sel:positive-marginality} implies that $\beta\ge\beta_{\mathrm m}$ and $q=q_+(\beta)$. This proves \eqref{sel:two-possible-plateaux}, and it shows that the second alternative is absent for $\beta<\beta_{\mathrm m}$. The formula for $\ell_\beta$ is \eqref{eq:plateau-multiplier} for the cluster $(c,r,\ell_\beta)$.
\end{proof}

\section{The dynamical transition and the selected energy}\label{sec:selected-sector}

Recall the dynamical transition parameters
\[
 \qd=\frac{p-2}{p-1},\qquad
 \bd^2=\frac{(p-1)^{p-1}}{p(p-2)^{p-2}}.
\]
We compare them with the constants in \eqref{sel:spinodal-constants}. First, $\qd-q_{\mathrm m}=(p-2)/(p(p-1))>0$. Second, a direct computation shows that $G(\qd)=p(p-2)^{p-2}/(p-1)^{p-1}=\bd^{-2}$. As $G$ decreases strictly on $[q_{\mathrm m},1]$, we find that $\bd^{-2}<G(q_{\mathrm m})=\beta_{\mathrm m}^{-2}$, that is, $\bd>\beta_{\mathrm m}$, and that $q_+(\bd)=\qd$. For $\beta>\bd$, we have $G(q_+(\beta))=\beta^{-2}<G(\qd)$, and then $q_+(\beta)>\qd$. For $\beta\ge\bd$, we write
\[
 q_\beta=q_+(\beta),
\]
so that $q_\beta\ge\qd$, with equality exactly at $\beta=\bd$.

\begin{theorem}[The selected stationary sector]\label{sel:selected-sector}
For every $p\ge3$ and $\beta\ge\bd$, the unique stationary plateau is $q_\beta$. In particular, the limit $\ell_\beta$ of $\lambda(t)$ and the limiting energy satisfy
\begin{align}
 \ell_\beta
 &=\frac1{1-q_\beta}+\beta^2p\bigl(1-q_\beta^{p-1}\bigr),
 \label{sel:selected-multiplier}\\
 e_\beta(\infty)
 &=\beta\bigl\{1+(p-2)q_\beta^{p-1}-(p-1)q_\beta^p\bigr\}.
 \label{sel:selected-energy}
\end{align}
For every $0\le A<\infty$,
\begin{equation}
 \sup_{0\le s\le t\le A}
 \bigl\{|C(T+t,T+s)-c(t-s)|+|R(T+t,T+s)-r(t-s)|\bigr\}
 \longrightarrow0\qquad(T\to\infty),
 \label{sel:stationary-window}
\end{equation}
where $r=-c'$ and $c(\tau)\to q_\beta$.
At the endpoint, $e_{\bd}(\infty)=\bd$ and $\ell_{\bd}=1+p\bd^2$.
\end{theorem}

\begin{proof}
In view of \cref{sel:all-temperature}, it suffices to exclude the alternative $q=0$. Suppose that $q=0$. Then \eqref{eq:plateau-multiplier} gives $\ell_\beta=1+\beta^2p$, which means that $\gamma=\ell_\beta-\beta^2p=1$ and
\[
 F(y)=\bigl(1+\beta^2py^{p-1}\bigr)(1-y)-1
 =y\bigl\{\beta^2py^{p-2}(1-y)-1\bigr\}.
\]
Since $\qd^{\,p-2}(1-\qd)=(p-2)^{p-2}/(p-1)^{p-1}$, we have $\beta^2p\,\qd^{\,p-2}(1-\qd)=\beta^2/\bd^2$ and
\[
 F(\qd)=\qd\left(\frac{\beta^2}{\bd^2}-1\right).
\]
For $\beta\ge\bd$, this is nonnegative. On the other hand, $\qd$ lies in $(0,1)=(q,1)$, so $F(\qd)<0$ by \eqref{eq:scalar-sign}. This contradiction shows that $q>0$.

At $\beta=\bd$, we have $F(\qd)=0$, and the contradiction comes from the strict inequality in \eqref{eq:scalar-sign}. In \cref{prop:scalar-clusters}, that strict inequality follows from the bound $c'(t)\ge F(c(t))$ and from $c'=-r<0$ at every finite time. In this way, the endpoint $\beta=\bd$ is included without any continuity argument in $\beta$.

We have shown that $q>0$, and \cref{sel:all-temperature} then implies that $q=q_+(\beta)=q_\beta$. The formula \eqref{sel:selected-multiplier} is the formula for $\ell_\beta$ in \cref{sel:all-temperature}, \eqref{sel:stationary-window} is \eqref{eq:all-temperature-window}, and the relation $r=-c'$ is part of \eqref{eq:cluster-fdt-properties}. For the energy, we use the marginality relation $\beta^2G(q_\beta)=1$. Multiplying it by $q_\beta/(p\beta(1-q_\beta))$, we obtain
\[
 \frac{q_\beta}{p\beta(1-q_\beta)}=\beta(p-1)q_\beta^{p-1}(1-q_\beta).
\]
Using $e_\beta(\infty)=(\ell_\beta-1)/(p\beta)$ and \eqref{sel:selected-multiplier}, we find
\[
 e_\beta(\infty)
 =\frac{q_\beta}{p\beta(1-q_\beta)}+\beta\bigl(1-q_\beta^{p-1}\bigr)
 =\beta(p-1)q_\beta^{p-1}(1-q_\beta)+\beta\bigl(1-q_\beta^{p-1}\bigr),
\]
and expanding the right side yields \eqref{sel:selected-energy}. At $\beta=\bd$, we have $q_\beta=\qd$ and $(p-1)\qd=p-2$, so that
\[
 (p-2)\qd^{\,p-1}-(p-1)\qd^{\,p}=\qd^{\,p-1}\bigl\{(p-2)-(p-1)\qd\bigr\}=0.
\]
This shows that $e_{\bd}(\infty)=\bd$, and \eqref{eq:energy-identity} gives $\ell_{\bd}=1+p\bd e_{\bd}(\infty)=1+p\bd^2$.
\end{proof}

The plateau and the energy in \cref{sel:selected-sector} are those predicted by Cugliandolo and Kurchan. In the normalization of \cite{ck1993}, the inverse temperature is $\sqrt2\,\beta$ and the energy is $-e/\sqrt2$; with this translation, \eqref{sel:selected-energy} is their Equation~(9), evaluated with the fluctuation--dissipation ratio $x=(p-2)(1-q)/q$ given after their Equation~(12).

\begin{remark}[The subsequent low-temperature limit]\label{rem:low-temperature}
As $\beta\to\infty$, we have $q_\beta\to1$, $\beta(1-q_\beta)\to[p(p-1)]^{-1/2}$, and
\[
 \lim_{\beta\to\infty}e_\beta(\infty)=2\sqrt{\frac{p-1}{p}}=\Eth.
\]
This agrees with the threshold energy of the pure $p$-spin landscape, computed by Crisanti and Sommers \cite[Equation~(9)]{cs1995} in the normalization of \cite{ck1993} and by Auffinger, Ben Arous, and \v{C}ern\'y \cite[Equation~(2.13)]{abc2013} in the present one. It also agrees with the results of Sellke \cite[Theorems~1.1 and~1.7]{sellke2024}, which show that, at low temperature, the energy of Langevin dynamics on a bounded time horizon comes close to $\Eth$. Here $\beta\to\infty$ is taken after the dimension limit and after both long-time limits, and the calculation says nothing about the gradient flow at zero temperature.

To verify these limits, recall that $q_\beta\ge\qd>0$. The marginality relation gives
\[
 (1-q_\beta)^2=\frac1{\beta^2p(p-1)q_\beta^{p-2}}\le\frac1{\beta^2p(p-1)\qd^{\,p-2}}\longrightarrow0,
\]
so $q_\beta\to1$. Taking square roots in the same relation, we find
\[
 \beta(1-q_\beta)=\frac{q_\beta^{-(p-2)/2}}{\sqrt{p(p-1)}}\longrightarrow\frac1{\sqrt{p(p-1)}}.
\]
Let $g(y)=1+(p-2)y^{p-1}-(p-1)y^p$, so that $e_\beta(\infty)=\beta g(q_\beta)$ by \eqref{sel:selected-energy}. We have $g(1)=0$ and
\[
 g'(y)=(p-1)y^{p-2}\bigl\{(p-2)-py\bigr\},\qquad g'(1)=-2(p-1).
\]
Let $\varepsilon=1-q_\beta$. By Taylor's theorem, there exists $\eta\in(q_\beta,1)$ such that
\[
 g(q_\beta)=g(1)-g'(1)\varepsilon+\tfrac12g''(\eta)\varepsilon^2
 =2(p-1)\varepsilon+\tfrac12g''(\eta)\varepsilon^2.
\]
On $[0,1]$, we have $|g''|\le K_p$ with $K_p=(p-1)\{(p-2)^2+p(p-1)\}$, and therefore
\[
 \bigl|e_\beta(\infty)-2(p-1)\beta(1-q_\beta)\bigr|\le\tfrac12K_p\,\beta(1-q_\beta)^2.
\]
As $\beta\to\infty$, the right side equals $\frac12K_p\cdot\beta(1-q_\beta)\cdot(1-q_\beta)$ and tends to zero, while $2(p-1)\beta(1-q_\beta)\to2(p-1)/\sqrt{p(p-1)}=\Eth$.
\end{remark}

\begin{remark}[The unselected intermediate interval]\label{sel:intermediate-range}
For $0<\beta<\beta_{\mathrm m}$, \cref{sel:all-temperature} gives $q=0$; then $\ell_\beta=1+\beta^2p$ and $e_\beta(\infty)=\beta$. For $\beta_{\mathrm m}\le\beta<\bd$, it proves convergence but allows either the zero branch or the positive branch $q_+(\beta)$. In the latter case, writing $q_+=q_+(\beta)$, we have
\[
 e_\beta(\infty)=\beta+
 \frac{q_+\{(p-2)-(p-1)q_+\}}{p\beta(p-1)(1-q_+)^2}>\beta.
\]
Indeed, the marginality relation can be written as $\beta^2pq_+^{p-1}=q_+/((p-1)(1-q_+)^2)$. Combining it with the formula for $\ell_\beta$ in \cref{sel:all-temperature}, we have
\[
 e_\beta(\infty)-\beta=\frac{\ell_\beta-1-\beta^2p}{p\beta}
 =\frac1{p\beta}\left\{\frac{q_+}{1-q_+}-\frac{q_+}{(p-1)(1-q_+)^2}\right\},
\]
and bringing the bracket to a common denominator gives the stated formula. For the sign, note that $\beta<\bd$ implies $G(q_+)=\beta^{-2}>G(\qd)$. As $G$ decreases strictly on $[q_{\mathrm m},1]$, we obtain $q_{\mathrm m}\le q_+<\qd$. In particular, $(p-2)-(p-1)q_+>0$, and the correction is positive. We do not decide which alternative occurs in this interval, and the transition is not classified from below there; see \cref{sec:discussion}.
\end{remark}

\section{Forgetting of fixed waiting times}\label{sec:forgetting}

\begin{theorem}[Forgetting of fixed waiting times]\label{sel:fixed-wait-forgetting}
Let $p\ge3$ and $\beta>0$, and let $(c,r)$ be the stationary profile of \cref{sel:all-temperature}. Then
\begin{equation}
 \lim_{t\to\infty}\sup_{0\le s\le S}
       \bigl\{C(t,s)+R(t,s)\bigr\}=0\qquad(0\le S<\infty).
 \label{sel:fixed-wait-decay}
\end{equation}
If the stationary plateau is zero, the stronger whole-row conclusion holds:
\begin{equation}
 \sup_{0\le s\le t}
 \bigl\{|C(t,s)-c(t-s)|+|R(t,s)-r(t-s)|\bigr\}
 \longrightarrow0\qquad(t\to\infty).
 \label{sel:zero-plateau-whole-row}
\end{equation}
\end{theorem}

\begin{proof}
The proof has four steps. We first reduce \eqref{sel:fixed-wait-decay} to a statement about $C$. When the plateau $q$ is positive, we verify the hypotheses of part~\ref{item:threshold-forgetting} of \cref{prop:threshold-inputs}. When $q=0$, we argue directly: the norm of the prefix center tends to zero when first $t\to\infty$ and then $L\to\infty$. Finally, we prove \eqref{sel:zero-plateau-whole-row} by splitting each row into a recent strip and a remote part.

\paragraph{Step 1: reduction to the correlation.}
Recall from \eqref{eq:response-derivative-bounds} that $0\le R(t,s)\le C(t,s)$ for $t\ge s\ge0$. For $t\ge S$, we have
\[
 \sup_{0\le s\le S}\bigl\{C(t,s)+R(t,s)\bigr\}\le2\sup_{0\le s\le S}C(t,s),
\]
and it suffices to show that $\sup_{0\le s\le S}C(t,s)\to0$ as $t\to\infty$.

\paragraph{Step 2: positive plateau.}
Suppose that $q>0$. Then the translated kernels converge along every sequence $T\to\infty$ to the single profile $(c,r)$ by \cref{sel:all-temperature}, and the plateau satisfies $\beta^2p(p-1)q^{p-2}(1-q)^2=1$ by \cref{sel:positive-marginality}. These are the hypotheses of part~\ref{item:threshold-forgetting} of \cref{prop:threshold-inputs}. Then \eqref{eq:threshold-input-forgetting} implies that $\sup_{0\le s\le S}|C(t,s)|\to0$.

\paragraph{Step 3: zero plateau.}
Suppose that $q=0$. Fix $L\ge0$ and let $t\ge L$. For $0\le s\le t-L$, the vector $v_s$ lies in $\mathcal H_s\subseteq\mathcal H_{t-L}$ by \cref{prop:geometry-tree-realization}, so $P_{t-L}v_s=v_s$. As $P_{t-L}$ is an orthogonal projection, we have
\[
 C(t,s)=\ip{v_t}{P_{t-L}v_s}=\ip{P_{t-L}v_t}{v_s}=\ip{M_{t-L}(L)}{v_s}.
\]
By the Cauchy--Schwarz inequality and $\norm{v_s}=1$, we obtain
\[
 \sup_{0\le s\le t-L}C(t,s)^2\le\norm{M_{t-L}(L)}^2.
\]
We now let $t\to\infty$ with $L$ fixed. By \cref{sel:all-temperature}, the translated kernels converge to the stationary cluster $(c,r,\ell_\beta)$ along every sequence $T_n\to\infty$. Then \cref{sel:prefix-reflection} applies along every such sequence, and \eqref{sel:prefix-norm} gives $\norm{M_{T_n}(L)}^2\to c(2L)$. As the sequence is arbitrary, we have $\norm{M_T(L)}^2\to c(2L)$ as $T\to\infty$. Setting $T=t-L$, we obtain
\begin{equation}
 \limsup_{t\to\infty}\sup_{0\le s\le t-L}C(t,s)^2\le c(2L).
 \label{sel:zero-plateau-remote-bound}
\end{equation}
This bound holds for every plateau, but it is useful only when $q=0$. Fix $0\le S<\infty$. For $t\ge S+L$, we have $[0,S]\subseteq[0,t-L]$, and \eqref{sel:zero-plateau-remote-bound} implies $\limsup_{t\to\infty}\sup_{0\le s\le S}C(t,s)^2\le c(2L)$ for every $L\ge0$. Letting $L\to\infty$ and using $c(2L)\to q=0$, we find that $\sup_{0\le s\le S}C(t,s)\to0$. Together with Steps~1 and~2, this proves \eqref{sel:fixed-wait-decay}.

\paragraph{Step 4: the whole row when $q=0$.}
Fix $L>0$, and let $t\ge L$. We split $[0,t]$ into the recent strip $[t-L,t]$ and the remote part $[0,t-L]$. On the recent strip, we use \eqref{eq:all-temperature-window} with $T=t-L$ and $A=L$. Writing $s=T+u$ with $u\in[0,L]$, we have $C(t,s)=C(T+L,T+u)$ and $c(t-s)=c(L-u)$, and similarly for $R$ and $r$. We obtain
\[
 \sup_{t-L\le s\le t}
 \bigl\{|C(t,s)-c(t-s)|+|R(t,s)-r(t-s)|\bigr\}\longrightarrow0
 \qquad(t\to\infty).
\]
On the remote part, we have $t-s\ge L$. Recall that $0\le R\le C$, that $0\le r\le c$, and that $c$ is nonincreasing. For $0\le s\le t-L$, we find
\[
\begin{aligned}
 |C(t,s)-c(t-s)|+|R(t,s)-r(t-s)|
 &\le C(t,s)+c(t-s)+R(t,s)+r(t-s)\\
 &\le2C(t,s)+2c(L).
\end{aligned}
\]
By \eqref{sel:zero-plateau-remote-bound}, the upper limit as $t\to\infty$ of the supremum of the right side over $0\le s\le t-L$ is at most $2c(2L)^{1/2}+2c(L)$. Combining the two parts, we have
\[
 \limsup_{t\to\infty}\sup_{0\le s\le t}
 \bigl\{|C(t,s)-c(t-s)|+|R(t,s)-r(t-s)|\bigr\}
 \le2c(2L)^{1/2}+2c(L),
\]
where $t\to\infty$ is taken with $L$ fixed. Since $q=0$, the right side tends to zero as $L\to\infty$. This proves \eqref{sel:zero-plateau-whole-row}.
\end{proof}

\section{Correlation aging}\label{sec:aging}

\begin{corollary}[Correlation aging in the ordered limits]\label{sel:aging}
For every $p\ge3$ and $\beta\ge\bd$,
\begin{equation}
 \lim_{\tau\to\infty}\lim_{s\to\infty}C(s+\tau,s)=q_\beta>0,
 \qquad
 \lim_{s\to\infty}\lim_{\tau\to\infty}C(s+\tau,s)=0.
 \label{sel:aging-ordered-limits}
\end{equation}
\end{corollary}

\begin{proof}
Fix $\tau\ge0$. Applying \eqref{sel:stationary-window} with $A=\tau$ and $T=s$ to the pair of offsets $(\tau,0)$, we find that $\lim_{s\to\infty}C(s+\tau,s)=c(\tau)$. As $\tau\to\infty$, we have $c(\tau)\to q_\beta$ by \cref{sel:selected-sector}, and $q_\beta\ge\qd>0$. This gives the first limit. Next, fix $s\ge0$. Applying \eqref{sel:fixed-wait-decay} with $S=s$ and $t=s+\tau$, we find that $C(s+\tau,s)\to0$ as $\tau\to\infty$. The inner limit is zero for every $s$, and so is the outer limit.
\end{proof}

The first inner limit in \eqref{sel:aging-ordered-limits} is the stationary correlation $c(\tau)$, and the second inner limit vanishes at every fixed waiting time. \Cref{sel:aging} concerns only these two ordered limits. It does not identify the time scale on which $C(s+\tau,s)$ decreases below $q_\beta$, and it does not give a fluctuation--dissipation relation in the slow regime. We discuss these questions in \cref{sec:discussion}.

\paragraph{Proof of \cref{thm:main-aging}.}
We collect the results of this section, which rely on \cref{prop:threshold-inputs}, proved in \cref{sec:tree-ceilings,app:general-selection}. Fix $p\ge3$ and $\beta>0$. By \cref{sel:all-temperature}, the multiplier and the energy converge to $\ell_\beta$ and $e_\beta(\infty)$, and the translated kernels converge locally uniformly to the stationary profile $(c_\beta,r_\beta)=(c,r)$. This is \eqref{eq:main-stationary}. In Step~3 of the proof of \cref{sel:all-temperature}, we showed that every stationary cluster equals $(c,r,\ell_\beta)$, so the profile is unique. By \cref{prop:stationary-clusters}, we have $c(0)=r(0)=1$ and $r=-c'$, and $c$ is completely monotone by \cref{sel:cluster-complete-monotonicity}. The forgetting statement \eqref{eq:main-forgetting} is \eqref{sel:fixed-wait-decay}. For $\beta\ge\bd$, \cref{sel:selected-sector} shows that the plateau is $q_\beta\ge\qd$, which is the larger solution of \eqref{eq:main-marginality}, and it gives \eqref{eq:main-energy} and \eqref{eq:main-multiplier}. Finally, \eqref{eq:main-aging} is \eqref{sel:aging-ordered-limits}. This completes the proof of \cref{thm:main-aging}.
